\chapter[Art: Imitation or Creation?]{Does Art Imitate the World or Create One?}
\section{A Painted Pipe}
First, pick up an object: an artwork hanging in a museum, apparently unremarkable, a pipe painted on canvas.

It is meticulously rendered. A polished brown wooden bowl, a black mouthpiece curving comfortably, exact proportions, reflections, shadows: almost a photograph from a tobacconist's window. Anyone approaching would say, "That is a pipe."

Beneath it, however, elegant French lettering reads "Ceci n'est pas une pipe": "This is not a pipe."

Belgian painter René Magritte made The Treachery of Images in 1929. First-time viewers nearly always pause, then smile: he is right. Tobacco cannot be packed into it; it cannot be lit or held between lips. It is paint on cloth. Confusion follows. If not a pipe, what is it? We recognized one immediately; how does recognition work? Why labor to make it resemble a pipe, then deny it in writing?

With a painting and a sentence, Magritte pinned an ancient question to the wall: does art imitate the world, or create a world?

If resemblance is art's skill, the sentence is nonsense. If the sentence is right and painting and pipe differ fundamentally, does resemblance still count as skill? Is the artist's brush a mirror or a wand?

Take the pipe with us as we begin.

\section{Tears in an Athenian Theater}
Travel back more than twenty-four centuries to Athens.

It is spring in the late fifth century BCE, the great Dionysian festival. The city takes a holiday, shops close, even prisoners receive temporary release for theater. The Theater of Dionysus rises along the southern slope in semicircular stone tiers seating over ten thousand. At dawn, people arrive with food and cushions: craftsmen in rough tunics, cloaked farmers, women holding children, nobles surrounded by servants. Olives, sheep's cheese, and humanity scent the air as dried-fig sellers weave between rows.

Today brings Sophocles' Oedipus the King. A signal quiets the crowd. Masked actors appear, thick-soled boots increasing their height. Plague afflicts Thebes; an oracle demands discovery of the former king's murderer. Oedipus vows to investigate. Witnesses assemble the truth: he himself killed the older king, his father, on a road years before, and married his own mother. The queen runs inside and kills herself; Oedipus takes golden pins from her dress and blinds himself.

Thousands follow the unfolding plot, fists clenched, breath caught. When the blinded king gropes out in bloodstained robes, many weep. Afterward they walk silently downhill, shaken for days.

Notice the oddity: nobody actually believes a murder occurred. The "blind" actor will remove mask and costume and drink at a tavern tonight. Everyone knows masks are plaster and linen, blood dye, palace scenery. Yet tears, fear, and the emptied-and-refilled feeling afterward are real.

Adults voluntarily spend money and an entire day watching things known to be false, becoming genuinely pained and shaken. That deserves explanation.

Among Athenians of every occupation sits a young man destined to wrestle with this forever: Plato. Tradition says the literary-minded young aristocrat wrote tragedies and dithyrambs himself. As a philosopher, one of his hardest questions would be this scene: why do unreal things have real power, and is that power blessing or danger?

\section{What Kind of Skill Is Resemblance?}
Lay out the conflict before answering.

Almost everyone has felt an intuition: artistic skill means making things look real.

Ancient and cross-cultural, it appears in Pliny the Elder's Natural History. Greek painter Zeuxis supposedly painted grapes so convincingly that birds pecked them. Rival Parrhasius invited him to inspect a painting. Zeuxis reached to uncover it, discovering the curtain itself was painted. One fooled birds; the other fooled the painter. Later generations treated deceiving the eye as painting's highest trophy.

China has a near-identical legend. Tang writer Zhang Yanyuan's Record of Famous Paintings through the Ages tells how Cao Buxing, painting a screen for Sun Quan in the Three Kingdoms era, turned an accidental ink spot into a fly. Sun Quan tried to flick it away. These may be legends rather than events, but legends reveal what people honor. "Just like the real thing" remains the easiest compliment for painting today.

An opposite intuition is equally persistent once considered:

If resemblance is everything, should cameras have unemployed every painter? A phone can take a hundred images in a second, each more realistic than a diligent painter. Why have phone shops not replaced museums? Edvard Munch's The Scream has a wax-melted face and irrationally red sky, failing a likeness test yet making countless viewers dizzy with recognized terror. Magritte's denial beneath his pipe is precisely its most famous feature.

The intuitions collide:

\textbf{First claim:} art imitates. The world exists; artists faithfully carry it into paint, words, or sculpture. More accurate imitation makes better art.

\textbf{Second claim:} art creates. The world is raw material for previously nonexistent colors, expressions, melodies, perspectives. Perfect imitation is merely photocopying.

This is not childish bickering. It affects whether art students begin with plaster-cast drawing, photographs' truth compared with paintings, abstract art as supposed emperor's clothes, learning to paint amid AI's many styles, and what prize judges actually judge.

Choose provisionally: is realistic likeness art's central skill? Remember your answer; the chapter will challenge it.

\section{Plato's Full Concern}
Plato first punctured the imitation problem. His famously severe ideal city politely expels poets. This readily becomes a joke about philosophers envying poets. Instead, \textbf{give Plato his full reasons}, strong enough for his recognition. Understanding a serious opponent is not agreeing.

His concern has four increasingly deep layers.

First, knowledge. Republic X distinguishes three beds: the Form, the carpenter's physical imitation, and the painter's imitation of that bed's appearance from one angle. Painting imitates imitation, three levels from truth. The painter need not know woodworking, weight-bearing legs, or joints, only appearances. Imitation captures the world's skin rather than its principles; resemblance may therefore be skill at departing from truth.

Second, psychology. Tragedy stimulates what everyday life should control: bereavement, fear, self-pity. A good person restrains grief in life but freely weeps with actors and enjoys it. Emotions resemble muscles: weekly indulgence strengthens them until they seize life's steering wheel. Fake scenes train real weakness.

Third, imitation itself shapes us. People do not simply watch; repeated performance and viewing make them become what they imitate. Villainous gestures and tones seep into actors, heroic sacrifice into spectators' courage. Imitation is a two-way pipe, carrying the stage into bodies. Precisely this formative power makes unrestricted art dangerous for the city.

Fourth, often missed, Plato does not hate beauty. His education includes music, poetry, gymnastics, beauty from childhood, but selectively: courage, moderation, harmony remain; cowardice, excess, wailing depart. He wants art as reason's assistant, not abolished. His enemy is \textbf{uncontrolled} beauty.

These reasons are not ridiculous. His real question is whether a power bypassing judgment to manipulate emotion needs supervision. Debates over video algorithms, violent games, and propaganda films still ask it.

Now his student Aristotle speaks for the defense.

The Poetics responds almost point by point. First, imitation is natural. Children learn through it, and people enjoy even accurate depictions of disgusting things because recognition itself pleases: "Ah, that is this." Imitation is not deception but a fundamental source of learning and pleasure.

Second, emotions are more than trouble to switch off. Aristotle's famous definition gives tragedy pity and fear leading to katharsis. Scholars have debated the word for two millennia: cleansing accumulated feelings, venting them safely, or refining them into appropriate form. Whatever the reading, theater becomes emotional exercise or purification rather than weakness fermenting. Rehearsing life's greatest disasters from safety may better equip you for life outside.

Third, poetry approaches philosophy more closely than history. History records particular past actions; poetry shows what a kind of person would do by probability or necessity. It deals in possibilities and human universals. A photograph shows this event; tragedy shows what human existence might encounter. Fiction is not falsehood and may contain deeper truth.

Teacher and student see opposite routes: art bypassing reason or reaching universality. Two millennia later, neither has decisively won.

A third voice sidesteps the dispute: "I simply paint what I see." Nineteenth-century French realist Gustave Courbet supposedly demanded to see an angel before painting one. Cameras seemed its final weapon: lenses have no position, negatives do not lie; machines can score perfectly at imitation.

Examine this separately, because it hides a crucial reminder: \textbf{realism is not absence of choice}.

Consider photography's strongest case. In 1936, Dorothea Lange photographed a mother and children at a California pea-pickers' camp. Migrant Mother became the Depression's emblem, seemingly reality entering a lens unassisted. Yet Lange took a series, approaching and varying angle and distance before choosing the published image. Choice began in noticing: whom to approach, when to shoot, which frame, which crop. A lens may not lie, but the person behind it judges every step.

Painters choose still more. Courbet's visible world required format, viewpoint, morning or evening light, inclusions and exclusions. Thousands of decisions precede faithful depiction. Realism is their result, not automatic copying. Even a casual meal photograph selects overhead or side view, filter or none, steaming food or scraps. \textbf{Resemblance is never simply the world's property; it is a result of choices.}

Three voices now stand together: dangerous imitation, useful imitation, undisputed recording. The third has immediately exposed its own limits. Perhaps the imitation-or-creation choice itself is mistaken.

\section[Four Questions for an Artwork]{A Bridge for Thought: Four Questions for an Artwork}
Can we analyze a painting, song, sculpture, or film without relying on talent or vague feeling? Two thousand years of aesthetics yield four portable questions. Ask each before an artwork.

\textbf{First: what does it represent? What does it resemble?}

Plato and Aristotle's central battlefield: apple, sea, face? How do perspective, light, proportion enable recognition? This trains attention, but is only one question. Magritte's sentence teaches that.

\textbf{Second: what does it make me feel? What does it express?}

Expression theory underlies this, famously formulated by Leo Tolstoy in What Is Art? A person communicates experienced feeling through sound, line, color, and words, infecting others with the same emotion. Art's core becomes sincere, strong transmission rather than resemblance. A painter first struck by dusk carries that impact through paint into you. The Scream exemplifies it: its distorted face resembles nobody, yet plants inexpressible terror in the viewer's chest. Under this question, the issue is contact, not likeness.

\textbf{Third: how is it organized? Is its form itself beautiful?}

Remove The Scream's colors and compare curved sky with straight bridge; suspend a poem's meanings and hear rhythmic pauses and surges. Formalism underlies this. Early-twentieth-century British critic Clive Bell called it "significant form": lines and colors arranged in particular relationships move us, not merely their depicted story.

Arts imitating little provide strong evidence. Mosques often avoid human and animal figures, yet lack no richness: endlessly extending geometry and Arabic calligraphy turn writing into breathtaking patterns. No world is directly copied, yet beauty remains in formal relations. Hokusai's Great Wave off Kanagawa likewise combines recognizable water and daring flat design: white claws of foam, Fuji's small triangle, expanses of Prussian blue. Its nineteenth-century arrival astonished European painters and changed Western painting. Representation and form work independently on one sheet.

Correct another widespread error. Nineteenth-century Europeans seeing elongated West African masks and geometric sculptures assumed primitive inability to draw people. Their distortion was often deliberate: enlarged foreheads signifying wisdom, simplified features invoking ancestors rather than individuals. Each departure had conceptual purpose. Picasso and other early modernists learned from such masks to break perspective and reconstruct images, opening modern art, while many objects had reached Europe through colonial looting. Nonlikeness often means another intention, not incapacity. \textbf{Reading deliberate nonlikeness as inability reveals the viewer's arrogance, not the maker's clumsiness.}

\textbf{Fourth: what does it say? What idea does it propose?}

Some works act primarily through ideas. Magritte's denial punctures habitual viewing with language. This question leads to our deepest waters, explored below.

Representation, expression, form, concept: four lights, not mutually expelling factions. One work can shine under all four, another under one especially. Dismissal for nonlikeness and praise for photographic accuracy each switch on only one lamp.

\section{Chinese Painters in Their Own Words}
Read primary materials. Chinese painters long spoke sharply about resemblance.

The Southern Dynasties' A New Account of Tales of the World, "Artistic Skill," recounts Eastern Jin painter Gu Kaizhi leaving figures' eyes unfinished for years. Asked why, he replied:

\begin{quote}
"Whether the limbs are beautiful or ugly has little to do with the essential marvel; conveying the spirit in a likeness lies precisely in these."
\end{quote}
The limbs' appearance matters less than the eyes for conveying a person's spirit; "these" translates the colloquial a-du.

Weigh the claim. Apparently a technique for eyes, it is really a program: portraiture seeks not accurate limbs but the elusive quality making someone themselves. Gu does not deny likeness but ranks it second. Form carries the cargo of spirit. Over sixteen centuries ago, Chinese painters clearly distinguished correct depiction from living depiction.

More than eight centuries later, Southern Dynasties writer Liu Xie discussed literary conception in The Literary Mind and the Carving of Dragons, "Imaginative Thought":

\begin{quote}
"Climb a mountain and feeling fills the mountain; behold the sea and thought overflows the sea."
\end{quote}
The direction matters: mountains do not pour shape into minds; human emotion flows outward across mountain and sea. Mind and object reverse imitation's direction. Art wants not a world awaiting copying but one steeped and colored in emotion. Liu expressed this more than a millennium before Tolstoy.

Do not turn these passages into the slogan "Chinese art expresses, Western art depicts." Medieval European icons prized sacred presence over accurate perspective, while Song court painters reached astonishing realism. Under Emperor Huizong, painters observed which foot peacocks lifted onto a mound and how rose stamens changed by hour, nearly scientific scrutiny. Civilizations are not sealed boxes with fixed personalities. Expressiveness and realism contend within each, Gu and Huizong differing as Plato and Aristotle did in Athens.

Primary materials are witnesses, not antiques. Gu and Liu testify that mere imitation never won uncontested in Chinese art theory; contemporary painters testify that likeness was never despised either. Multiple voices make actual history.

\section[Cave Art: Several Explanations]{Deep in the Caves: One Fact, Several Explanations}
Go deeper to humanity's earliest art and practice competing explanations of shared facts.

French and Spanish caves, Lascaux, Altamira, Chauvet, preserve paintings ten, twenty, even over thirty thousand years old: Lascaux roughly seventeen thousand, Chauvet perhaps beyond thirty thousand. Animals dominate: bison, horses, reindeer, mammoths. Their quality was so high that experts initially declared Altamira's 1879 discoveries forged, refusing to believe Paleolithic people could paint so well. Later work overturned doubt. Muscles, fur, movement reveal sustained precise observation. Establish this: \textbf{these were not clumsy primitives; their eyes and hands would merit professional standing today}.

Yet subjects receive unequal treatment. Animals live vividly; humans often become stick figures. Realism as choice was already operating twenty millennia ago. Painters deliberately selected what deserved detail.

Why paint, and there? Most works lie not at mouths but in dark depths reached by torch, narrow passages, and hundreds of meters of travel, without daylight or residents. This is not living-room decoration. Skilled work, deep locations, detailed animals, schematic humans: several explanations compete.

\textbf{First: hunting magic.} Late-nineteenth- and early-twentieth-century scholars including Abbé Breuil argued that depicting prey magically controlled it. A drawn bison was virtually captured; arrowlike marks and repeatedly overlaid animals supported the idea. This reconnects art to subsistence, perhaps a life-or-death technology. But dietary bone remains often differ from painted animals: horses and bison painted, reindeer eaten. Handprints and abstract signs also resist the hunting-blueprint account.

\textbf{Second: art for art's sake.} Perhaps people simply liked beauty and admiration. This modest account avoids pretending to know minds. Yet why hide pictures in hundreds of meters of darkness rather than inhabited, firelit rock shelters with audiences?

\textbf{Third: ritual and shamanism.} Late-twentieth-century accounts associated with David Lewis-Williams connect caves with religious experience. Darkness, reduced oxygen, and flickering light promote altered consciousness; cave walls become a membrane between reality and spirit, animals summoned from beyond. This explains location as ritual space, animals as spirits, abstract signs as geometric visions. Its weakness is scarce direct evidence for minds thirty thousand years ago and heavy cross-cultural analogy, risking a uniform shamanic mold for diverse societies.

None defeats the others outright. That is a lesson: \textbf{explanation is a searchlight whose disciplinary position, economics, aesthetics, religion, illuminates one face and leaves others dark}. No living witness tells us what painters understood. Humility is essential. One piece of evidence plus one story is not possession of truth.

All three agree more deeply: accurate likeness was not the ultimate purpose. It was a means to another end, twenty thousand years ago as now.

\section[Why Can a Urinal Be Art?]{A Deeper Challenge: Why Can a Urinal Be Art?}
Enter the deepest waters, where art history's famous prank meets serious theory.

New York, 1917. The Society of Independent Artists promised an exhibition without juries: pay the fee and exhibit. Marcel Duchamp submitted a shop-bought ceramic urinal, sideways, signed "R. Mutt" in black paint and titled Fountain.

Organizers were stunned. After argument, it was excluded: a supposedly threshold-free exhibition discovered its threshold. Duchamp resigned. A small magazine then published a defense through another writer's pen, broadly: whether Mutt made it mattered little; he \textbf{chose} it. A new title, angle, and setting removed utility and created a new idea for an everyday thing.

He claimed neither resemblance, since a urinal resembles itself, nor beauty, though some admire its curve. Instead, \textbf{art may designate rather than manufacture, reframing rather than representing the world}. Of our four questions, representation and expression scarcely apply, form only slightly. Concept does the real work.

The original disappeared; museum examples are replicas authorized in Duchamp's lifetime. Its question remains: if a signed urinal in a museum is art, where does non-art begin?

Philosophers took it up. In 1964, Arthur Danto saw Andy Warhol's Brillo Boxes, wooden boxes looking identical to supermarket soap cartons, and wrote his famous artworld account. Some objects count as art not through visible difference but through a world of theory, art-historical knowledge, and interpretive atmosphere supporting them. Appearance alone cannot decide.

George Dickie developed a more explicit institutional theory: artworks are artifacts granted candidate status for appreciation within artworld practices involving museums, critics, artists, and audiences. Art status depends partly on institutional recognition, not objects alone. Deflating perhaps, but familiar: paper becomes money through institutions; land becomes a border through agreement. Art may be similarly constituted.

The theory explains Fountain as conceptual interrogation of art and an identical shop urinal as outside that arena. It also accommodates conceptual and performance art where resemblance ceases to help.

Three limits matter. It answers art status, not quality: institutional acceptance does not establish profundity; judgment still needs looking, feeling, arguing. It risks insider control: who authorizes the artworld, and can a clique merely promote itself? And it can eclipse craft and emotion. A child's unrecognized, earnest drawing may fail its definition, yet does that concentration really differ from Gu painting eyes? The bright line marks institutional identity while leaving humanity's ancient artistic impulse outside its light.

All four lamps are now lit. Art can \textbf{represent} the world, as Aristotle defends; \textbf{express} feeling, as Tolstoy and Liu explain; \textbf{organize form}, as Bell and Islamic patterns show; and \textbf{propose concepts}, as Duchamp and Danto explore. Not rival monopolies on definition, but different human activities often combined in masterpieces. The either-or dissolves: imitation creates through selection; creation borrows the world's lines and colors even in abstraction.

\section{Filters, Beautification, and AI Artists}
The old question happens daily in your pocket.

A selfie already contains algorithmic selections: smoothing, slimming, brightening, adjustable parameters. Supposedly recording your real self, you create a little each time. Filters are contemporary styles, comparable to Courbet's framing choices with buttons moved into software. Plato might point and say we still bathe in imitations of imitations.

Consider photography's similar shock in 1839. A painter supposedly declared painting dead. It survived, but faced a hard question: once machines mastered likeness, what remained for brushes? Impressionists explored light and fleeting sensation, Post-Impressionists structure and emotion, abstraction discarded likeness. As photography took representation, painting explored expression, form, and concept. Each copying technology forces humanity to reconsider art's remaining work.

Now AI generates convincing images from a line of text in seconds and supplies styles on request. From Zeuxis's grapes to Lange's camera to diffusion models, resemblance's price approaches zero. Our questions help: AI is remarkably strong at representation and formal combination, while expression, experiencing and transmitting an impact, and concept, what the work questions, still require human answers. This remains contested, not settled; some hold that genuine transmission matters regardless of human or machine origin. Your generation inherits the argument.

\section{Your Question: Likeness or Something Else?}
Return to the pipe.

Magritte's "not a pipe" forces recognition that pictures are not the world itself; art stands somewhere outside. What does it do there? Two great voices still contend.

One demands witness: hunger, war, an expression need accurate preservation. Art without roots in reality evades responsibility; a painter abandoning reality betrays the pea-picking mother. The other demands imagination: records are plentiful, new ways of seeing, unnamed feelings, and possible alternatives scarce. Art's value creates what does not yet exist. If it only repeats the world while not being the world, what is it?

Give reasons for whether art matters more as accurate depiction or creation of the previously nonexistent. There is no standard answer. But at your next photograph, painting, or AI image, you have four lamps: what it represents, makes you feel, organizes, and says. Which lamps you switch on already answers the chapter's question.
